%==============================================================
% ΚΕΦΑΛΑΙΟ 6 — ΣΥΜΠΕΡΑΣΜΑΤΑ ΚΑΙ ΜΕΛΛΟΝΤΙΚΕΣ ΕΠΕΚΤΑΣΕΙΣ
%==============================================================
\chapter{Συμπεράσματα και Μελλοντικές Επεκτάσεις}
\label{ch:conclusions}

Η παρούσα εργασία εξέτασε την επικύρωση, ερμηνεία και παρακολούθηση
μοντέλων μηχανικής μάθησης στον πιστωτικό κίνδυνο, με έμφαση στο πώς
μεταβάλλονται οι απαιτήσεις αυτές καθώς αυξάνεται η πολυπλοκότητα του
μοντέλου. Στο παρόν κεφάλαιο συνοψίζονται τα ευρήματα και διατυπώνονται
οι προτεινόμενες κατευθύνσεις για περαιτέρω έρευνα.

%==============================================================
\section{Σύνοψη της Εργασίας}
\label{sec:summary}
%==============================================================

Η εργασία ξεκίνησε με ανασκόπηση του θεωρητικού και κανονιστικού
πλαισίου του πιστωτικού κινδύνου (Κεφάλαιο~\ref{ch:theory}) και
συνέχισε με συστηματική σύγκριση τριών μοντέλων στη βάση FICO HELOC
(Κεφάλαια~\ref{ch:methodology} και~\ref{ch:results}): ενός Logistic
Regression ως παραδοσιακού σημείου αναφοράς, ενός XGBoost ως κύριου
μοντέλου ανάλυσης και ενός Feed-Forward Neural Network ως
επιδεικτικού benchmark πολυπλοκότητας. Η αξιολόγηση πραγματοποιήθηκε
σε τέσσερις άξονες ---απόδοση, ερμηνευσιμότητα, μεροληψία και
παρακολούθηση--- και τα ευρήματα συντέθηκαν σε μια συγκεκριμένη
πρόταση πλαισίου επικύρωσης (Κεφάλαιο~\ref{ch:framework}).

%==============================================================
\section{Απάντηση στο Κεντρικό Ερευνητικό Ερώτημα}
\label{sec:research_question}
%==============================================================

Το κεντρικό ερώτημα της εργασίας ήταν πώς επικυρώνονται, ερμηνεύονται
και παρακολουθούνται τα μοντέλα μηχανικής μάθησης στον πιστωτικό
κίνδυνο και πώς μεταβάλλονται οι απαιτήσεις όταν αυξάνεται η
πολυπλοκότητα του μοντέλου. Η απάντηση που προκύπτει είναι ότι η
αύξηση της πολυπλοκότητας δεν συνεπάγεται γραμμική αύξηση των
απαιτήσεων επικύρωσης. Ορισμένα στάδια ---όπως η παρακολούθηση
σταθερότητας μέσω PSI και CSI ή η ανίχνευση μεροληψίας μέσω DPD και
DI--- παραμένουν ουσιαστικά αμετάβλητα ανεξάρτητα από το μοντέλο.
Αντίθετα, η ερμηνευσιμότητα είναι ο άξονας στον οποίο η πολυπλοκότητα
έχει δραματική επίπτωση: η μετάβαση από τη γραμμικότητα της
Logistic Regression στη μη γραμμική δομή του XGBoost εισάγει μέτριο
επιπλέον κόστος μέσω του TreeExplainer, ενώ η περαιτέρω μετάβαση στο
FFNN πολλαπλασιάζει το κόστος αυτό κατά τάξεις μεγέθους χωρίς
αντίστοιχο όφελος στην απόδοση. Η ορθή απάντηση στο ερώτημα
επικύρωσης ενός σύνθετου μοντέλου δεν είναι λοιπόν μια ομοιόμορφη
κλιμάκωση όλων των ελέγχων, αλλά μια στοχευμένη ενίσχυση στους
άξονες όπου η πολυπλοκότητα προκαλεί πραγματική απώλεια διαφάνειας.

%==============================================================
\section{Κύρια Συμπεράσματα}
\label{sec:main_conclusions}
%==============================================================

Από τη συστηματική ανάλυση των τεσσάρων αξόνων αξιολόγησης προκύπτουν
τέσσερα κύρια συμπεράσματα.

\paragraph{Απόδοση.}
Τα τρία μοντέλα παρουσίασαν οριακές διαφορές στους δείκτες AUC, Gini
και KS, με το XGBoost να προηγείται ελαφρώς. Το εύρημα αυτό
οφείλεται κυρίως στην εκτεταμένη προεπεξεργασία των χαρακτηριστικών
από την FICO, η οποία γραμμικοποιεί τις σχέσεις μεταξύ τους και
περιορίζει το performance gap. Συνεπάγεται ότι η επιλογή μοντέλου
στον πιστωτικό κίνδυνο δεν μπορεί να γίνεται αποκλειστικά με κριτήριο
την απόδοση, καθώς αυτή ενδέχεται να είναι παρόμοια μεταξύ μοντέλων
διαφορετικής πολυπλοκότητας.

\paragraph{Ερμηνευσιμότητα.}
Ο άξονας αυτός αποτελεί το βασικό σημείο διαφοροποίησης μεταξύ των
μοντέλων. Το XGBoost, μέσω του TreeExplainer, παρέχει ολοκληρωμένη
SHAP ανάλυση σε δευτερόλεπτα, ικανοποιώντας πλήρως το κανονιστικό
«δικαίωμα εξήγησης». Αντίθετα, το FFNN απαιτεί KernelExplainer με
υπολογιστικό κόστος της τάξης των ωρών, καθιστώντας την ερμηνεία
πρακτικά απαγορευτική σε παραγωγικά περιβάλλοντα με μεγάλο όγκο
δεδομένων. Το γεγονός αυτό αποτελεί ισχυρό επιχείρημα κατά της
χρήσης βαθιών νευρωνικών δικτύων στον πιστωτικό κίνδυνο όταν δεν
υπάρχει αντίστοιχο όφελος απόδοσης.

\paragraph{Μεροληψία.}
Και τα τρία μοντέλα εμφάνισαν σχεδόν ταυτόσημες τιμές DPD
(${\sim}0{,}33$) και DI (${\sim}0{,}47$), σαφώς εκτός των αποδεκτών
ορίων. Η σύμπτωση αυτή αποδεικνύει ότι η πηγή της μεροληψίας ήταν τα
δεδομένα και όχι η εσωτερική λογική κάποιου μοντέλου. Επιπλέον, ο
μετριασμός μεροληψίας δεν είναι δωρεάν: η in-processing προσέγγιση
στη Logistic Regression κόστισε ${\sim}14\%$ της AUC, ενώ η
post-processing προσέγγιση στο XGBoost διατήρησε την AUC αλλά
μείωσε το Balanced Accuracy. Η επιλογή τεχνικής μετριασμού αποτελεί
επιχειρηματική απόφαση και πρέπει να γίνεται με ρητή αξιολόγηση του
σχετικού κόστους.

\paragraph{Παρακολούθηση.}
Οι δείκτες PSI και CSI λειτούργησαν αξιόπιστα ως σύστημα έγκαιρης
προειδοποίησης. Στα τρία σενάρια οικονομικής επιδείνωσης, ο PSI
παρήγαγε σταδιακά αυξανόμενες τιμές
($0{,}069 \to 0{,}210 \to 0{,}641$ για το XGBoost), επιτρέποντας
παρέμβαση πριν η απόδοση υποβαθμιστεί. Σημαντικό εύρημα αποτελεί ότι
η ίδια η παρακολούθηση είναι ανεξάρτητη της πολυπλοκότητας του
μοντέλου, καθώς οι δείκτες αυτοί υπολογίζονται επί των δεδομένων ή
των προβλέψεων χωρίς να εξαρτώνται από την εσωτερική δομή του
αλγορίθμου. Η κρίσιμη διαφορά εμφανίζεται κατά τη \emph{διάγνωση}
ενός drift event: εδώ το χαμηλό κόστος ερμηνείας του XGBoost γίνεται
αποφασιστικό πλεονέκτημα έναντι του FFNN.

Συνολικά, τα τέσσερα συμπεράσματα ενισχύουν την κεντρική θέση της
εργασίας ότι το XGBoost αποτελεί τη βέλτιστη επιλογή στο σημερινό
πλαίσιο του πιστωτικού κινδύνου, καθώς συνδυάζει επαρκή απόδοση,
αποδοτική ερμηνευσιμότητα και ευελιξία στον μετριασμό μεροληψίας
χωρίς τις απαγορευτικές απαιτήσεις των βαθιών νευρωνικών δικτύων.

%==============================================================
\section{Μελλοντικές Επεκτάσεις}
\label{sec:future_work}
%==============================================================

Οι περιορισμοί που αναγνωρίστηκαν στην ενότητα~\ref{sec:limitations}
ορίζουν φυσικές κατευθύνσεις για περαιτέρω έρευνα.

\paragraph{Επέκταση σε πολλαπλά datasets.}
Η εφαρμογή του προτεινόμενου πλαισίου σε διαφορετικές βάσεις
δεδομένων ---όπως τα Lending Club, German Credit ή ευρωπαϊκά datasets
καταναλωτικής πίστης--- θα επέτρεπε αξιόπιστο έλεγχο της
γενικευσιμότητας. Datasets χωρίς εκτεταμένη προεπεξεργασία
χαρακτηριστικών αναμένεται να εμφανίσουν μεγαλύτερο performance gap
μεταξύ των μοντέλων, μετατοπίζοντας πιθανώς την ισορροπία μεταξύ
απόδοσης και ερμηνευσιμότητας.

\paragraph{Ανάλυση μεροληψίας σε πραγματικά δημογραφικά χαρακτηριστικά.}
Η χρήση της \texttt{AverageMInFile} ως proxy αποτελεί μεθοδολογική
αναγκαιότητα λόγω της φύσης του HELOC. Η εφαρμογή του πλαισίου σε
βάσεις με πραγματικά δεδομένα φύλου, ηλικίας ή εθνικότητας θα
παρείχε ουσιαστικότερη επικύρωση των μηχανισμών ανίχνευσης και
μετριασμού, σε άμεση συμφωνία με τις απαιτήσεις των ρυθμιστικών
αρχών.

\paragraph{Δυναμική παρακολούθηση με χρονικές σειρές.}
Η αντικατάσταση της στατικής προσομοίωσης drift από πραγματικά
χρονικά δεδομένα θα επέτρεπε τη μελέτη συσχετισμένων μετατοπίσεων σε
πολλαπλά χαρακτηριστικά ταυτόχρονα και τη βαθμονόμηση των ορίων PSI
και CSI σε πραγματικές οικονομικές συνθήκες. Πιθανές πηγές δεδομένων
αποτελούν χρηματοοικονομικές βάσεις που καλύπτουν περιόδους κρίσης
(όπως 2008--2009 ή 2020).

\paragraph{Επέκταση σε LGD και EAD.}
Η παρούσα εργασία επικεντρώθηκε στην πρόβλεψη της PD. Η αντίστοιχη
ανάλυση των παραμέτρων LGD και EAD ---οι οποίες ολοκληρώνουν τον
υπολογισμό της Αναμενόμενης Ζημίας κατά Basel III--- θα προσέφερε
πληρέστερο πλαίσιο επικύρωσης για ολόκληρο τον κύκλο διαχείρισης
του πιστωτικού κινδύνου.

\paragraph{Σύγκριση εναλλακτικών αλγορίθμων Gradient Boosting.}
Η αντιπαραβολή του XGBoost με τις υλοποιήσεις LightGBM και CatBoost
στο ίδιο πλαίσιο θα μπορούσε να αναδείξει αν η υπεροχή του XGBoost
στις απαιτήσεις επικύρωσης διατηρείται ή αν κάποια από τις
εναλλακτικές προσφέρει καλύτερη ισορροπία ταχύτητας, ακρίβειας και
ερμηνευσιμότητας. Παρόμοια συγκριτική ανάλυση θα μπορούσε να γίνει
και μεταξύ των explainers (TreeExplainer, KernelExplainer,
DeepExplainer) ως προς το trade-off ακρίβειας έναντι υπολογιστικού
κόστους.

\paragraph{Αυτοματοποιημένα συστήματα ειδοποίησης.}
Τέλος, η ενσωμάτωση των δεικτών παρακολούθησης σε μηχανισμούς
αυτόματης ειδοποίησης ---με προκαθορισμένα κατώφλια και αυτοματοποιημένη
επανεκπαίδευση--- αποτελεί φυσική συνέχεια του πλαισίου προς ένα
ολοκληρωμένο MLOps περιβάλλον για τον πιστωτικό κίνδυνο, σε ευθεία
συμφωνία με τη φιλοσοφία της συνεχούς επικύρωσης που προβλέπεται
από τις EBA Guidelines.